\section{Mean jets and the phase of the end contributions}\label{sec:phase}
The action comparison in Proposition~\ref{arr:action-upper} may take place
at complex common parameters.  It is therefore necessary to determine the
end phases to $o(1)$ inside the exponential in Proposition~\ref{jump:jump}.
Polynomial closeness of the end profiles by itself would not suffice.
We first compute the mean opening jets, then compare a finite Taylor
graph with an actual stationary solve before the pole.  Integrating the
mean equation transfers this information to the height matching section
and determines the exponent through its constant term.  This precision
is what allows contributions from several ends to be compared.
Throughout this section an opening scale is $B=|s|^{-k/2}$, and
$|x|=O(B^{-2})$.  The Taylor orders and derivative orders invoked below,
together with their polynomial losses, are fixed before $B$ tends to
infinity.  Where derivatives of actual solutions are used, the indicated
strip and first-trace norms retain their scale factors.

On a chosen cylindrical end, use the physical radial coordinate after
the analytic axis adjustments of Proposition~\ref{lin:family}.  The
additional first-order tilt is removed as in
Proposition~\ref{arr:preliminary-solves}.  Write
\begin{equation}
 w=R^2/2-1,\qquad t=\log z,\qquad
 \langle f\rangle=(2\pi)^{-1}\int_0^{2\pi}f(\theta)\,d\theta.
 \label{phase:w}
\end{equation}
The opening coefficient $\beta=\beta(b,x)$ is the first-jet coefficient;
it is linear homogeneous in $x$ for fixed $b$, and real analytic in $b$.

\begin{lemma}[Mean jets and their improved growth]\label{phase:jets}
Let $w_{[n]}$ denote the term homogeneous of degree $n$ in $x$ in the
stationary jet, with $b$ retained as an analytic parameter.  For each
fixed $n$ and every sufficiently small $\eta>0$, the estimates below
hold uniformly for $b$ in a smaller complex neighborhood, with fixed
logarithmic and angular derivatives:
\begin{align}
 \langle w_{[0]}\rangle&=O(z^{-4+\eta}),\notag\\
 \langle w_{[1]}\rangle
       &=\beta(z^2-2)+O(|x|z^{-2+\eta}),\notag\\
 \langle w_{[2]}\rangle
       &=(-4\beta^2\log z+c_2(b,x))z^2
                                  +O(|x|^2z^{1+\eta}),
 \label{phase:mean-jets}\\
 w_{[n]}&=O_n(|x|^n z^{2n-2+\eta})\qquad(n\ge2).
 \label{phase:all-jets}
\end{align}
Here $c_2$ is quadratic homogeneous in $x$, analytic in $b$, and real
for real parameters.
\end{lemma}
\begin{proof}
We use the full radial equation, so that its angular and axial terms can
be counted at every opening degree.  Put
\[
 A=1+w,\qquad H=A+\frac{w_\theta^2}{4A},\qquad f=\log A.
\]
Here $A$ and $H$ are analytic units near the cylindrical base profile,
and $f$ is the logarithm continued from $A=1$.  To derive the equation, set
$\pi=w_z$, $\chi=w_\theta$, and $D=2H+\pi^2$.  The radial curvature numerator
preceding \eqref{lin:cylinder-operator}, together with the support term
$(R-zR_z)/2$, is the radial shrinker equation.  Multiplying it by $RD$
and collecting its terms gives
\[
 \begin{aligned}
 0={}&-H\left(w_t-2w-\frac{4A w_{\theta\theta}-4\chi^2}{4A^2+\chi^2}\right)
       +2H w_{zz}-\frac{\pi\chi}{A}w_{z\theta}\\
 &+\pi^2\left(w-1-\frac{w_t}{2}+\frac{w_{\theta\theta}}{2A}\right).
 \end{aligned}
\]
Here $D$ is the determinant factor of the radial graph metric.  Since
\begin{equation}
 \frac{4A w_{\theta\theta}-4w_\theta^2}{4A^2+w_\theta^2}
       =2\partial_\theta\arctan(f_\theta/2),
 \label{phase:angular-divergence}
\end{equation}
division by $H$ and conversion to $t=\log z$ yield the exact equation
\begin{equation}
 \begin{split}
 w_t-2w={}&2\partial_\theta\arctan(f_\theta/2)
       +2z^{-2}(w_{tt}-w_t)
       -z^{-2}\frac{w_t w_\theta}{AH}w_{t\theta}\\
 &+z^{-2}\frac{w_t^2}{H}
       \left(w-1-\frac{w_t}{2}+\frac{w_{\theta\theta}}{2A}\right).
 \end{split}
 \label{phase:full-equation}
\end{equation}
The arctangent is the analytic branch near zero.  In particular its
angular derivative has zero mean at every formal degree.  When $w$ is
circular, \eqref{phase:full-equation} is equivalently
\begin{equation}
 w_{zz}-\frac z2w_z+w
       =\frac{w_z^2}{4(1+w)}(2+zw_z-2w).
 \label{phase:circular}
\end{equation}
The prescribed axis adjustments are ambient orthogonal transformations
before the physical radius is used, so these identities remain valid for
their parameter-dependent formal expansions.

Let $\mathfrak E(w)$ denote the left side minus the right side of
\eqref{phase:full-equation}, and put
$\mathcal D_b=D\mathfrak E(w_{[0]})$.  Outside the observation support,
$\mathfrak E(w_{[0]})=0$.  Proposition~\ref{lin:family} and the
analytic change from $R$ to $w$ give, after removal of the first-order
tilt,
\begin{equation}
 w_{[0]}=O(z^{-2+\delta}),\qquad
 w_{[1]}=\beta z^2+O(|x|z^\delta)
 \label{phase:initial-jets}
\end{equation}
for every small $\delta>0$, with each fixed log and angular derivative.
The leading first-order term is circular.  Each $O(z^\alpha)$ remainder
here and in the induction below is bounded in the weighted strip domain
$X_{z^\alpha,h}$ after each required fixed log derivative, together
with its first traces.  By taking $h$ large these give the displayed
pointwise angular bounds.  An unscaled log derivative used as a
coefficient is the function trace of a differentiated jet, supplied by
the log-regularity assertion of Lemma~\ref{lin:rates}.

The operator $\mathcal D_b$ is the radial realization of the Jacobi
operator at $w_{[0]}$: the preceding multipliers and change of dependent
variable are units, and their derivatives multiplying the base equation
vanish.  At the round cylinder one has exactly
\[
 D\mathfrak E(0)v
 =v_t-2z^{-2}(v_{tt}-v_t)-v_{\theta\theta}-2v
 =\mathcal L_{\mathrm{cyl}}v.
\]
For $w_{[0]}$ satisfying \eqref{phase:initial-jets}, the coefficients
of $v_{\theta\theta},v_\theta,v_t,v$ differ from these model coefficients
by $O(z^{-2+\delta})$ in the normalized multiplier norms.  The
$v_{tt}$ coefficient is unchanged, and the $v_{t\theta}$ coefficient
is $O(z^{-6+\delta})$.  Thus it is also small after division by the
mixed-derivative scale $z^{-1}$.  The same statements hold after each
fixed log derivative: they follow by differentiating the rational
coefficients in \eqref{phase:full-equation} and using
\eqref{phase:initial-jets}.  On a sufficiently far tail, uniformly for
$b$ in a smaller complex neighborhood, $\mathcal D_b$ therefore has
the normalized decay and strict diffusion margins of
Lemma~\ref{lin:rates}.

For $n\ge2$, the equation at total opening degree $n$ is
\begin{equation}
 \mathcal D_b w_{[n]}
 =-\left[\mathfrak E\left(w_{[0]}+
                    \sum_{1\le j<n}w_{[j]}\right)\right]_{[n]}.
 \label{phase:jet-equation}
\end{equation}
In particular every occurrence of the unknown $w_{[n]}$, including
the terms with nonconstant base coefficients, stays on the left.
The right side has only lower-degree jets.  To estimate it, first
consider the angular term in the rational form
\eqref{phase:angular-divergence}.  A Taylor coefficient of its
numerator contains an angular derivative of a jet.  A base angular
derivative is $O(z^{-2+\delta})$; an angular derivative of $w_{[1]}$
is $O(|x|z^\delta)$ because $\beta z^2$ is circular; and at degree
$j\ge2$ the inductive bound is $O(|x|^jz^{2j-2+\delta})$.
Each is two powers below the nominal weight $z^{2j}$ for its degree,
with degree zero interpreted as nominal weight one.  The inverse
denominator has degree-$\ell$ coefficients of at most the nominal weight
$|x|^\ell z^{2\ell}$ after a small loss.  Hence, for a requested final loss
$\eta$, the input losses can be chosen so that this angular contribution
at degree $n$ has weight at most $|x|^n z^{2n-2+\eta/3}$.

Every remaining nonlinear term in \eqref{phase:full-equation} has the
explicit factor $z^{-2}$.  Its rational coefficients are analytic in
$w,w_\theta$, and its other factors are fixed log or angular derivatives
of the lower jets.  At degree $n$ their nominal weight before this
factor is at most $|x|^n z^{2n+\eta/3}$ after the same choice of input
losses.  For the displayed second derivatives use bulk bounds as follows:
$w_{\theta\theta}$ uses its ordinary bulk slot, $w_{t\theta}$ uses
the function bulk norm of $\partial_t w$ at one extra angular Sobolev
order, and $w_{tt}$ uses the function bulk norm of $\partial_t^2w$.
These unscaled bounds come from the separately controlled differentiated
jets.  The other factors use their first traces.  This gives the same
source weight $|x|^n z^{2n-2+\eta/3}$ in
$Y_{z^{2n-2+\eta/3},h}$.  There are finitely many products at each
degree, so the input losses can always be chosen this small.

For clarity about derivative order, fix a terminal degree $N$ and a
required number $m$ of log and angular derivatives.  At degree $j\le N$
carry $m+2(N-j)$ derivatives, with a fixed angular Sobolev reserve.
The two derivatives in \eqref{phase:full-equation} then require only
estimates already obtained at lower degrees.  Apply the separated
inverse of Lemma~\ref{lin:rates} to \eqref{phase:jet-equation} at the
fixed weight $2n-2+2\eta/3>2$, taking the entrance sufficiently far out
for this weight.  For real $b$, solve at this weight with
the entrance data of the polynomially growing jet from
Proposition~\ref{lin:family}.  Compare in a larger polynomial weight
chosen above both this weight and the known polynomial bound for that
jet.  Both weights exceed every slow rate and prescribe the same full
slow entrance data, so separated uniqueness identifies the two
solutions.  This also excludes the fast outward mode; the slow
homogeneous $z^2$ mode lies in the target weight.

The target-weight inverse, source, and entrance trace are analytic in
$b$ on a smaller complex neighborhood.  The constructed solution is
therefore analytic there.  On each compact strip the original finite
jet is also analytic in $b$ by Proposition~\ref{lin:family}; equality
for real $b$ extends to that complex neighborhood by the identity
principle.  Thus the weighted estimate is uniform there without using
a real comparison principle at complex parameters.  Reserve the remaining
$\eta/3$ for the arbitrarily small loss in the fixed log-regularity
assertion.  It gives all the derivative bounds just specified at weight
$2n-2+\eta$.  This proves
\eqref{phase:all-jets} for each fixed $n$, using no convergence of the
formal series.

We now project the full equation onto its mean.  Define
\[
 \begin{gathered}
 \mathcal L_0=\partial_z^2-(z/2)\partial_z+1,\qquad
 U=(AH)^{-1},\quad V=H^{-1},\\
 K=w-1-\frac{w_t}{2}+\frac{w_{\theta\theta}}{2A}.
 \end{gathered}
\]
Averaging \eqref{phase:full-equation} cancels the angular divergence
exactly and gives
\begin{equation}
 \mathcal L_0\langle w\rangle
 =\frac1{2z^2}\left\langle
          U w_t w_\theta w_{t\theta}-V K w_t^2\right\rangle.
 \label{phase:mean-equation}
\end{equation}
The following coefficient bounds make the first three degrees explicit.
Use smaller losses in the input estimates and relabel their finite sums
by $\delta$; then choose $\delta$ small enough that the products below
have total loss at most the final $\eta$.  The bounds
\eqref{phase:initial-jets} and the now proved degree-two case of
\eqref{phase:all-jets} give
\[
 \begin{array}{c|cc}
 n&(w_t)_{[n]}&(w_\theta,w_{t\theta},w_{\theta\theta})_{[n]}\\ \hline
 0&O(z^{-2+\delta})&O(z^{-2+\delta})\\
 1&2\beta z^2+O(|x|z^\delta)&O(|x|z^\delta)\\
 2&O(|x|^2z^{2+\delta})&O(|x|^2z^{2+\delta}).
 \end{array}
\]
Here a bound for a tuple applies to each entry, with the fixed derivative
norms above.  Expanding the three rational factors through degree two
gives
\[
 \begin{array}{c|ccc}
       &[0]&[1]&[2]\\ \hline
 U&1+O(z^{-2+\delta})&O(|x|z^2)&O(|x|^2z^4)\\
 V&1+O(z^{-2+\delta})&-\beta z^2+O(|x|z^\delta)&O(|x|^2z^4)\\
 K&-1+O(z^{-2+\delta})&O(|x|z^\delta)&O(|x|^2z^{2+\delta}).
 \end{array}
\]
For example,
\[
 K_{[1]}=w_{[1]}-\tfrac12(w_{[1]})_t
       +\frac{(w_{[1]})_{\theta\theta}}{2A_{[0]}}
       -\frac{(w_{[0]})_{\theta\theta}w_{[1]}}{2A_{[0]}^2}.
\]
The circular terms in its first two summands cancel.  At degree two,
the possible factor $w_{[1]}^2$ in an inverse denominator is
$O(|x|^2z^4)$, and whenever it occurs in $K_{[2]}$ it multiplies
$(w_{[0]})_{\theta\theta}=O(z^{-2+\delta})$.  This proves the stated
$K_{[2]}$ bound and accounts for the denominator terms as well.

Set $\Xi=U w_t w_\theta w_{t\theta}$ and $\Psi=V K w_t^2$ for this
calculation.  Taking every distribution of degrees among the displayed
factors gives
\begin{equation}
 \begin{aligned}
 \Xi_{[0]}&=O(z^{-6+\eta}),&\Psi_{[0]}&=O(z^{-4+\eta}),\\
 \Xi_{[1]}&=O(|x|z^{-2+\eta}),&\Psi_{[1]}&=O(|x|z^\eta),\\
 \Xi_{[2]}&=O(|x|^2z^\eta),&
 \Psi_{[2]}&=-4\beta^2z^4+O(|x|^2z^{2+\eta}).
 \end{aligned}
 \label{phase:mean-product-bounds}
\end{equation}
To see the only leading term in the last entry, use the exhaustive
product decomposition
\[
 \Psi_{[2]}=(w_t^2)_{[2]}(VK)_{[0]}
          +(w_t^2)_{[1]}(VK)_{[1]}
          +(w_t^2)_{[0]}(VK)_{[2]}.
\]
The first factor in its first summand is
$4\beta^2z^4+O(|x|^2z^{2+\eta})$, and $(VK)_{[0]}=-1+O(z^{-2+\delta})$.
The other two summands are respectively $O(|x|^2z^{2+\eta})$ and
$O(|x|^2z^\eta)$.  In $\Xi_{[2]}$, either a positive-degree angular
factor loses two powers, or both angular factors have their base decay;
the coefficient bounds above give the displayed estimate for every
degree distribution.  The terms linear in $w_{[2]}$ through decaying
base coefficients are included here only after its growth bound has
been proved.  Thus \eqref{phase:mean-equation} contains no other term
at the resonant weight, and yields
\begin{equation}
 \begin{aligned}
 \mathcal L_0\langle w_{[0]}\rangle&=O(z^{-6+\eta}),\\
 \mathcal L_0\langle w_{[1]}\rangle&=O(|x|z^{-2+\eta}),\\
 \mathcal L_0\langle w_{[2]}\rangle
             &=2\beta^2z^2+O(|x|^2z^\eta).
 \end{aligned}
 \label{phase:mean-forcing}
\end{equation}
The loss in these source estimates is arbitrary.  In applying the scalar
estimate below, take it smaller than the $\eta$ in the statement and
reserve the difference for the fixed log-regularity estimate.

It remains to extract the slow mean coefficient.  Let $h_0(z)=z^2-2$,
so that $\mathcal L_0h_0=0$.  For a polynomially growing mean solution
$h_0v$ on a far tail, direct differentiation gives
\[
 (h_0^2e^{-z^2/4}v')'=h_0e^{-z^2/4}\mathcal L_0(h_0v).
\]
Polynomial growth excludes the fast homogeneous solution.  Hence
\begin{equation}
 v'(z)=-\frac{e^{z^2/4}}{h_0(z)^2}
          \int_z^\infty h_0(s)e^{-s^2/4}\mathcal L_0(h_0v)(s)\,ds.
 \label{phase:slow-flux}
\end{equation}
For a source $O(z^\sigma)$ this is $O(z^{\sigma-3})$ by the Gaussian
tail estimate.  When $\sigma<2$, integration shows that the solution is
$c h_0+O(z^\sigma)$; if the solution is $o(z^2)$, then $c=0$.
The scalar separated estimate gives the same bounds with the fixed log
derivatives when the source has them.

Apply this observation to the first line of \eqref{phase:mean-forcing},
relaxing its source weight to $-4+\eta$.  The base bound in
\eqref{phase:initial-jets} excludes $h_0$, and gives the asserted
$O(z^{-4+\eta})$ mean.  For degree one subtract $\beta h_0$.
The remaining function is $O(|x|z^\delta)=o(z^2)$ by
\eqref{phase:initial-jets}, and its source has weight $-2+\eta$.
This proves the first two lines of \eqref{phase:mean-jets}.

Finally,
\[
 \mathcal L_0(-4\beta^2z^2\log z)
       =2\beta^2z^2-8\beta^2\log z-12\beta^2.
\]
After subtracting this particular solution, the last source in
\eqref{phase:mean-forcing} is $O(|x|^2z^{1+\eta})$ for any fixed
$0<\eta<1$, allowing a harmless relaxation of its lower weight.
Apply \eqref{phase:slow-flux} with
\[
 v_2(z)=\frac{\langle w_{[2]}\rangle+4\beta^2z^2\log z}{h_0(z)}.
\]
It gives $v_2'=O(|x|^2z^{-2+\eta})$, which is integrable.  Thus
$c_2=\lim_{z\to\infty}v_2(z)$ exists and gives
\[
 \langle w_{[2]}\rangle+4\beta^2z^2\log z
       =c_2(z^2-2)+O(|x|^2z^{1+\eta}).
\]
The coefficient is $O(|x|^2)$ and homogeneous quadratic in $x$,
because both the finite jet and the subtracted term have that degree.
Absorbing $-2c_2$ into the remainder proves the third line of
\eqref{phase:mean-jets}.  The flux integral is uniformly convergent
on a smaller complex $b$ neighborhood, and the finite-order jets there
are analytic in $b$.  Evaluation at one fixed height together with that
integral proves that $c_2$ is analytic.  The same formula uses only real
data for real parameters, so $c_2$ is real there.
\end{proof}

\begin{lemma}[Evaluation of the inner mean]\label{phase:inner-evaluation}
Fix a sufficiently small $\epsilon_1>0$, and put
$z_1=B^{1-\epsilon_1}$ on the real incoming portion of an adverse end.
For the preliminary normalized solution at the common parameters,
\begin{equation}
 \frac{\langle w(z_1)\rangle}{z_1^2}
 =\beta-\frac{2\beta}{z_1^2}-4\beta^2\log z_1+c_2(b,x)
                                                   +o(B^{-4}).
 \label{phase:inner-value}
\end{equation}
The error is uniform over the finite list of ends and the allowed
parameter balls.
\end{lemma}
\begin{proof}
Put $Z=z_2=B^{1-\epsilon_1/2}$, and stop every cylindrical end at a
comparable real incoming section, retaining the compact core and the
conical domains.  Fix a Taylor degree $n$ and a small loss $\eta>0$.
Since $|x|z_2^2=O(B^{-\epsilon_1})$, the finite Taylor surface stays in
the same small graph charts on this domain.  Its residual, in the
normalized strip source norm at height $z$, is bounded by
\[
       C_n B^{-2n-2}(1+z)^{2n+2+\eta}.
\]
Its compact and conical residuals have the same coefficientwise Taylor
bound.  Let $e$ be the difference from the actual preliminary solution,
using the analytic scalar variable $w$ on the incoming cylinders.
The actual mean-value linearization, with its compact multiplier columns
and observation equations retained, gives
\[
 \mathcal A_{Z,n}(e,\ell)=(f_n,0),\qquad
 \mathcal A_{Z,n}(v,m)=(L_*v+Q_*m,P_*v).
\]
Here $P_*e=0$ follows from the identical observations, and does not set
$\ell$ to zero.  If a finite chart conversion leaves an observation
residual of Taylor order $n+1$, first remove it by a fixed compact lift;
its contribution has the same bound as $f_n$.

First solve the forced augmented problem with zero observations and
homogeneous outgoing derivative data.  Retain the natural Taylor weight
\[
       \rho_n(z)\asymp(1+z)^{p_n},\qquad p_n=2n+2+\eta,
\]
normalized on a fixed entrance collar.  Apply
Proposition~\ref{lin:global} at this fixed rate with its size parameter
relabeled $Z=z_2$, rather than the opening scale $B$.  The straight
real stops are at uniformly comparable multiples of $Z$; choose the
proposition's $c$ below their minimum ratio.  There are no modified
outer bands or height charts in this application.  The paths retain
the single real log coordinate, and the truncation leaves the compact
domain and infinite conical pieces unchanged.  Since
$|\beta|Z^2=O(B^{-\epsilon_1})$,
the circular coefficients retain their strict smallness and diffusion
margins on the truncated cylinders.  The normalized domain-to-range
smallness of the actual mean-value coefficients is the estimate
$Z(W_{\mathrm{act}}+W_{\mathrm T})=o(1)$ verified below solely from
the preliminary and finite-jet bounds for the eventual choice
$n\epsilon_1>2$.  The opening-dependent costs decrease on earlier
strips, while the type-preserving base tail is made small by the fixed
entrance choice.  Thus the truncated paths satisfy the
coefficient and compatibility hypotheses of the proposition.

For each finite conical stop, the index construction in the proof of
Proposition~\ref{arr:preliminary-solves} applies directly on this shorter
domain.  Its normalized principal tensor still has positive real part,
and the new derivative exit in the real log coordinate is transverse.
The same
fixed-domain interpolation to $-\Delta+1$ is therefore properly elliptic
with complementing boundary conditions, so its index is zero.
For this fixed $n$, choose the entrance sufficiently far out that
$p_n\sup E$ is small, and then take $Z$ large.  The exterior forward
estimate retains its strict buffer above every slow rate, and its
compact-limit argument still has polynomial growth, hence lies in the
Gaussian augmented domain.  The stopped index argument supplies
surjectivity.  The original $B$ enters the resulting bound through
the source amplitude $B^{-2n-2}$, so the forced pair obeys
\[
 \|(e_F,\ell_F)\|_{X_{\rho_n}}\le C_n B^{-2n-2},\qquad
 \frac{|e_F(z_1)|}{z_1^2}
 \le C_nB^{-2-2n\epsilon_1+(1-\epsilon_1)\eta}.
\]
The second estimate is the function trace of the retained weighted norm.
No coarse physical inverse factor $B^{q_n}$ is spent on this polynomial
source; increasing $p_n$ would have to be charged to the displayed loss.

The remainder $h=e-e_F$ has homogeneous augmented equation and zero
observations, but has outgoing derivative discrepancy $g$ at each stop.
Its norm in the natural Neumann trace space
\eqref{lin:neumann-trace} is at most $B^{M_n}$.
Construct a separate collar lift: choose fixed $0<a<b<c<1$, and solve
\[
 L_*H=0\quad\hbox{on }[az_2,z_2],\qquad
 H(az_2)=0,\qquad H_t(z_2)=g.
\]
This zero entrance condition is imposed on $H$, not on $h$.
The collar has bounded logarithmic length and $E\asymp z_2^{-2}$.
The shifted smallness hypothesis can be checked independently of this
comparison.  Use, on the whole real collar $z\asymp Z$,
the same circular reference
\[
 R_c(z)=\sqrt{2(1+\beta z^2)},\qquad w_c=\beta z^2.
\]
Because $|\beta|Z^2=O(B^{-\epsilon_1})$, this reference has uniform
transversality, nonzero drift and strict cylindrical diffusion margins,
with $E\asymp Z^{-2}$.  This verifies its admissibility; the correction
must satisfy a separate estimate.  The comparison construction in
Proposition~\ref{arr:preliminary-solves} and
\eqref{arr:preliminary-bound} give, in the full norm
\eqref{gauge:norm},
\[
 W_{\mathrm{act}}\le C_D\bigl(
 Z^{-2+\eta}+B^{-2}Z^\eta+B^{-4}Z^{2.2}\bigr).
\]
The first two terms are the base cylindrical remainder and the nonleading
first opening jet after the common axis and tilt adjustments; the last
is the actual correction to that comparison.  These estimates do not
use the Taylor comparison or the shielding conclusion.

For the degree-$n$ Taylor graph, Lemma~\ref{phase:jets}, including its
fixed logarithmic and angular derivative bounds, gives
\[
 W_{\mathrm{T}}\le C_n\bigl(
 Z^{-2+\eta}+B^{-2}Z^\eta+B^{-4}Z^{2+\eta}
 +(B^{-2}Z^2)^{n+1}\bigr).
\]
Indeed, for $j\ge2$ the corresponding term is bounded by
$C_nB^{-4}Z^{2+\eta}(B^{-2}Z^2)^{j-2}$, and there are only finitely
many terms.  The final summand can be omitted if the finite jet is taken
in $w$ and the radius is recovered exactly.  If the radius itself is
Taylor-truncated, it bounds the degree-$n+1$ tail of the analytic square
root; the same bound holds after each required fixed log derivative.
The analytic change between $R-R_c$ and $w-w_c$ preserves these strip
bounds, since $R_c$ stays bounded away from zero and its fixed log
derivatives are bounded.

Writing $\alpha=1-\epsilon_1/2$, the two estimates imply
\begin{align*}
 Z(W_{\mathrm{act}}+W_{\mathrm{T}})\le C_n\bigl(&
 B^{-\alpha(1-\eta)}
 +B^{-1-\epsilon_1/2+\alpha\eta}
 +B^{-0.8-1.6\epsilon_1}\\
 &+B^{-1-3\epsilon_1/2+\alpha\eta}
 +B^{1-(n+3/2)\epsilon_1}\bigr)=o(1).
\end{align*}
Here the existing choice $n\epsilon_1>2$, followed by sufficiently
small $\eta>0$, makes every exponent negative.  Interpolate the actual
and Taylor graphs in the common scalar chart used by the mean-value
linearization.  Their difference from this same circular reference is
bounded throughout the segment by a constant times
$W_{\mathrm{act}}+W_{\mathrm{T}}$.  In the radial realization of
Lemma~\ref{gauge:mixed} take $c_1=0$, $c_2=1$, and axial scale $Z$.
Thus the shifted coefficient--module bound of Lemma~\ref{lin:shielding}
applies uniformly along the averaged segment, with
$\varepsilon=O(Z(W_{\mathrm{act}}+W_{\mathrm{T}}))=o(1)$.


Lemma~\ref{lin:shielding}, with large scale $z_2$, gives on
$[bz_2,cz_2]$ the strip and first-trace bound
$B^{M'_n}e^{-c_0z_2^2}$.
Its application uses the shifted coefficient--module norm of the actual
difference operator; a coefficient containing a background second
derivative is kept in its $L^2$ slot.  Under the constant shift
$d=c_1z_2^2$, the exit is $(\partial_t+d)\widetilde H=g$.
Keeping this Robin condition and $E(\partial_t+d)^2$ in the shifted
norm preserves the fast-root estimate and the gain $d-O(Ed^2)$.

Choose $\chi$ with bounded log derivatives, zero below $bz_2$ and one
above $cz_2$.  Extend $V=\chi H$ by zero inward, summing over the finite
list of ends.  Its observations vanish, it realizes every remote datum,
and its residual is already exponentially small:
\[
       r=L_*V=[L_*,\chi]H,\qquad
       \|r\|_Y\le B^{M''_n}e^{-c_0z_2^2}.
\]
Only the shielded lift has been cut off; an arbitrary polynomially
bounded remote discrepancy would not give this estimate.
Solve the normalized repair
$\mathcal A_{Z,n}(k,\ell_R)=(-r,0)$ with homogeneous exits.
For this solve rate $2.2$ and any fixed polynomial conversion loss are
harmless, giving
\[
        |k(z_1)|/z_1^2+|\ell_R|
                \le B^{C_n}e^{-c_0z_2^2}.
\]
Since $z_1/z_2\to0$, $V$ vanishes near $z_1$.
The pair $(e,\ell)-(e_F,\ell_F)-(V+k,\ell_R)$ has zero augmented
output, zero observations and homogeneous exits in the same stopped
cylindrical and infinite conical domain.  Full augmented injectivity
therefore makes it zero.  This excludes the remaining slow mode together
with any compensating compact multiplier.

Combining the two solves and writing $\eta'=(1-\epsilon_1)\eta$ yields
\[
 \frac{|e(z_1)|}{z_1^2}
 \le C_n B^{-2-2n\epsilon_1+\eta'}
          +B^{C_n}e^{-c_0B^{2-\epsilon_1}}=o(B^{-4})
\]
when $n\epsilon_1>2$ and the losses are sufficiently small.
The associated trace statement controls
$(e,e_\theta,\sqrt E\,e_t)/z_1^2$; an unscaled derivative trace may
cost $E^{-1/2}$.  The later primitive estimate instead uses the separate
preliminary bounds $\|\delta_t\|_{L^2}\le CW$ and
$\|\delta_t\|_{\mathrm{trace}}\le CBW$, so it needs no stronger
Taylor-error derivative trace.

For $3\le j\le n$, \eqref{phase:all-jets} gives
\[
 B^{-2j}z_1^{2j-4+\eta}
       =B^{-4-(2j-4)\epsilon_1+(1-\epsilon_1)\eta}=o(B^{-4}).
\]
The divided base and linear errors are respectively
$O(z_1^{-6+\eta})$ and $O(B^{-2}z_1^{-4+\eta})$; the divided
second-order remainder is $O(B^{-4}z_1^{-1+\eta})$.
All are $o(B^{-4})$ for the chosen losses.  The surviving terms in
\eqref{phase:mean-jets} therefore give \eqref{phase:inner-value}.
Only finitely many Taylor orders have been used.
\end{proof}

\begin{lemma}[The integrated averaged equation]\label{phase:primitive}
Continue the preliminary solution from $z_1$ to its fixed real-radius
height section $r_0\in(0,\sqrt2)$.  In prescribed transition gauges,
use the circular averages of the scaled axial and unscaled radial
positions.  Denote the resulting mean curve again by $(z,R)$ and set
$w=R^2/2-1$.  Then
\begin{equation}
 \frac{d}{d\log z}\left(\frac w{z^2}\right)
       =\frac{4\beta}{z^2}-\frac{4\beta^2}{1+\beta z^2}
           +\mathcal E,
 \qquad \int_{z_1}^{z(r_0)}\mathcal E\,d\log z=o(B^{-4}).
 \label{phase:integrated-equation}
\end{equation}
\end{lemma}
\begin{proof}
The preliminary profile bound, with a slight relaxation of its powers,
gives size
\begin{equation}
 W=O(B^{-1.7})
 \label{phase:W}
\end{equation}
for the deviation from the circular model on each strip in this region.
Choose $\epsilon_1$ small enough for this consequence of
Proposition~\ref{arr:preliminary-solves}; the number of strips is
$O(\log B)$.  On the real radial portion $E\asymp z^{-2}$; on the
bounded final transition bands $E\asymp B^{-2}$.
Their derivative losses are consequently at most $B$ and $B^2$.
Lemma~\ref{gauge:mixed}, applied directly on the prescribed real parameter
bands, and \eqref{gauge:averaging} show that the mean curve obeys the
circular equation with a strip $L^2$ error at most
$C(BW^2+B^2W^3)$.  This accounts for the full angular curvature and
support terms, including their dependence on axial speed.
The axial coordinate of the mean has nonzero scaled derivative; conversion
of this circular curve to its circular radial equation uses bounded
factors.  It does not invert a nonconstant unknown complex coordinate.

For explicit normalization, the radial-velocity equation of a circular
curve multiplied by $R(1+R_z^2)$ is the left side minus the right side
of \eqref{phase:circular}.  This multiplier is bounded here.  Dividing
that equation to obtain the derivative of $w/z^2$ costs $O(z^{-2})$.
The accumulated noncircular error is therefore at most
\begin{align}
 C z_1^{-2}(BW^2+B^2W^3)(1+\log B)
 &\le C\{B^{-4.4+2\epsilon_1}+B^{-5.1+2\epsilon_1}\}(1+\log B)
 \notag\\[-2pt]&=o(B^{-4}).
 \label{phase:tensor-budget}
\end{align}
The difference between the mean of $R^2/2-1$ and the squared mean radius
is $\langle(R-\langle R\rangle)^2\rangle/2=O(W^2)$.
At either endpoint, after division by $z^2$, its cost is
$O(W^2 z_1^{-2})=o(B^{-4})$.  At the initial radial section it relates
the mean in Lemma~\ref{phase:inner-evaluation} to the present mean curve;
at the final pure height section the radius is exactly $r_0$.

It remains to control the circular substitution, including the second
derivative.  Put $\tau=\log z$ and $q=\beta z^2$.  The exact circular
identity is
\begin{equation}
 \partial_\tau(w/z^2)
 =2z^{-4}(w_{\tau\tau}-w_\tau)
 -\frac{w_\tau^2}{2z^4(1+w)}(2+w_\tau-2w).
 \label{phase:exact-primitive}
\end{equation}
Its value at $w=q$ is exactly the first two terms of
\eqref{phase:integrated-equation}.
For $\delta=w-q$, the strip and trace bounds give
\[
 \|\delta\|_{L^\infty}+\|\delta_\tau\|_{L^2}\le CW,
       \qquad \|\delta_\tau\|_{\mathrm{trace}}\le CBW.
\]
The denominator $1+q$ stays uniformly away from zero up to the chosen
section.  Since $BW\to0$, the rational expression in the second term
of \eqref{phase:exact-primitive} is Lipschitz on the permitted traces;
use one $L^2$ factor for $\delta_\tau$ in its products.  Its integrated
substitution error is bounded by
$Cz_1^{-4}W(1+\log B)=o(B^{-4})$.
For the first term integrate by parts once, obtaining the exact formula
\begin{equation}
 \int 2z^{-4}(\delta_{\tau\tau}-\delta_\tau)\,d\tau
       =[2z^{-4}\delta_\tau]_{\mathrm{ends}}
                   +6\int z^{-4}\delta_\tau\,d\tau.
 \label{phase:ibp}
\end{equation}
It costs at most
\begin{equation}
 C BWz_1^{-4}+C Wz_1^{-4}(1+\log B)
              =O(B^{-4.7+4\epsilon_1})+o(B^{-4}).
 \label{phase:second-budget}
\end{equation}
One may take $0<\epsilon_1<0.1$ and then the jet losses sufficiently
small.  Both \eqref{phase:tensor-budget} and
\eqref{phase:second-budget} are then $o(B^{-4})$.
This proves the integrated statement with all changes of gauge included.
\end{proof}

\begin{proposition}[The phase to constant order]\label{phase:expansion}
For an adverse end set $Z_0=(-\beta)^{-1/2}$ on the branch continued from
positive values, and define $Z$ by its mean height at the common section
as in Proposition~\ref{jump:jump}.  Let
\[
 \phi(r)=\sqrt{1-r^2/2},\qquad
 (1/r-r/2)d_0'+d_0/2=-\phi''/\phi'^2,\qquad d_0(r_0)=0,
\]
and write $d_0=2\log r+d_*+O(r^2(1+|\log r|))$.
With
\[
 Z_*=Z+Z^{-1}\{d_*-2\log(Z/2)\},
\]
one has
\begin{equation}
 Z_*^2=-\frac1\beta-8\log Z_0+\frac{c_2(b,x)}{\beta^2}
                                  +C(r_0)+o(1),
 \label{phase:expansion-eq}
\end{equation}
where $C(r_0)=2+2\log2$ for the stated normalization $d_0(r_0)=0$
and definition of $Z_*$, independently of $r_0$.  The logarithmic
coefficient $-8$ is the same for every end, and the error is uniform
on the finite set of adverse ends.  All logarithms retain their
prescribed continuation from the common incoming contours.
\end{proposition}
\begin{proof}
Write $q_0=\phi(r_0)^2\in(0,1)$.  At the mean section,
$z=Z\sqrt{q_0}$ and $w=-q_0$.
Integrate Lemma~\ref{phase:primitive} starting with
\eqref{phase:inner-value}.  An antiderivative of its explicit right side is
\[
 -2\beta/z^2-4\beta^2\log z+2\beta^2\log(1+\beta z^2).
\]
At $z_1$, the last logarithm is $o(1)$, since
$\beta z_1^2=O(B^{-2\epsilon_1})$.
Consequently
\begin{equation}
 -\frac{q_0}{z^2}=\beta-\frac{2\beta}{z^2}+c_2
       -4\beta^2\log z+2\beta^2\log(1+\beta z^2)+o(B^{-4}).
 \label{phase:section-equation}
\end{equation}
The preliminary estimate gives $Z/Z_0=1+O(B^{-1.8})$.
Equation~\eqref{phase:section-equation} first improves this to
$Z^2-Z_0^2=O(1+\log B)$: its left side plus $-\beta$ has size
$O(B^{-4}(1+\log B))$, and its derivative with respect to $Z^2$ has size
comparable to $B^{-4}$.  Thus all subsequent Taylor errors are $o(B^{-4})$
in that equation.  Writing $Z^2=-1/\beta+\Delta$, expansion gives
\begin{equation}
 \Delta=-4\log Z_0+\frac{c_2}{\beta^2}
       +\frac2{q_0}-2\log q_0+2\log(1-q_0)+o(1).
 \label{phase:Z-phase}
\end{equation}
All the displayed constants apart from $c_2/\beta^2$ are real.

For completeness the equation for $d_0$ simplifies exactly to
\begin{equation}
           (d_0/\phi)'=\frac2{r\phi^4}.
 \label{phase:d0}
\end{equation}
Its real initial value at $r_0$ therefore gives a real $d_*$ and the
stated logarithmic expansion.  Squaring the definition of $Z_*$ adds
$2d_*-4\log(Z/2)+o(1)$; the square of the correction is
$O(B^{-2}(1+\log B)^2)=o(1)$.
Since $\log Z-\log Z_0=o(1)$, combining this with
\eqref{phase:Z-phase} gives \eqref{phase:expansion-eq}, with
\[
 C(r_0)=\frac2{q_0}-2\log q_0+2\log(1-q_0)+2d_*+4\log2.
\]
The matching-radius constant can also be evaluated explicitly.
For $q=1-r^2/2$ on the positive real interval, put
\[
       F(q)=\frac1q-\log q+\log(1-q).
\]
Since $dq=-r\,dr$ and $r^2=2(1-q)$,
\[
 \frac{2\,dr}{r\phi(r)^4}
       =-\frac{dq}{q^2(1-q)}=dF(q).
\]
Consequently the initial condition $d_0(r_0)=0$ gives
\[
 d_0(r)=\sqrt q\,[F(q)-F(q_0)],\qquad q_0=1-r_0^2/2.
\]
As $r\downarrow0$, $F(q)=2\log r+1-\log2+O(r^2)$, whence
\[
 d_*=1-\log2-F(q_0)
       =1-\log2-\frac1{q_0}+\log q_0-\log(1-q_0).
\]
Thus the constant in \eqref{phase:expansion-eq} is independent of $r_0$:
\[
 C(r_0)=2F(q_0)+2d_*+4\log2=2+2\log2.
\]
All logarithms here have positive real arguments.  Their prescribed
continuation agrees with the incoming branches already used in the
phase calculation, so this simplification introduces no additional phase.

The chosen branches agree by continuation along the common incoming
part of the contours.  No $2\pi i$ ambiguity is inserted when an end
is compared with another having the same leading opening.
\end{proof}

\begin{theorem}[Exclusion of a critical opening arc]\label{phase:no-arc}
Every real formal arc $(b(s),x(s))$ through the origin
satisfying $\nabla\mathcal F(b(s),x(s))=0$, where $\mathcal F$ is
the formal Gaussian action of \Cref{lin:family}, has $x\equiv0$.
\end{theorem}
\begin{proof}
Suppose instead that
$x=s^k v+O(s^{k+1})$, $v\ne0$.
The invertible first-jet opening map has at least one nonzero leading
coefficient
\[
 \beta_j=c_js^k+O(s^{k+1}),\qquad c_j\in\mathbb R.
\]
If a $c_j$ is negative, start at the positive $s$ ray.
If all nonzero $c_j$ are positive, start at the ray
$\arg s=\pi/k$.  This choice also applies when $k$ is even:
it is a continuation on the chosen sheet of $s$, and does not
require replacing a real parameter by its negative.
Lemma~\ref{arr:contours} and Proposition~\ref{arr:action-upper} give both
lateral arrays and common parameters in an open range of nearby test
angles.  Their common parameters have the same real formal jet, to any
prescribed finite order, as the formal implicit substitution used there.
The two continuations need not be complex conjugates.

We make precise the reality information retained by actual substitution.
For any fixed required order, choose the polynomial center and implicit
orders high enough that the actual common parameters agree with their
real formal expansion to an error of that order.  To expand $-1/\beta_j$
through its constant term it suffices to know $\beta_j$ through order
$2k$, with remainder $o(s^{2k})$.  The array expansion provides
this and more.  If $c_{2,j}(b,x)=Q_j(x)+o(|x|^2)$ at $b=0$, then on
any of the fixed complex rays
\begin{equation}
 \frac{c_{2,j}(b,x)}{\beta_j^2}
     \longrightarrow \frac{Q_j(v)}{c_j^2}\in\mathbb R.
 \label{phase:real-ratio}
\end{equation}
The identical factor $s^{2k}$ cancels.  In particular, leading
reality of the opening after ray rotation is not being used to assert
reality of its unknown subleading values.

Let $P_j(s)$ be the negative-power principal part of the Laurent
series $-1/\beta_j$.  All its coefficients, including the separate
constant coefficient of that Laurent series, are real.
For each pair with $P_i\ne P_j$, let $as^{-m}$ be the first
nonzero term of their difference.  On
$s=\rho e^{i\vartheta}$ its real part is
$a\rho^{-m}\cos(m\vartheta)$.  Only finitely many angles in a small
interval are excluded by its vanishing.  Choose one fixed angle avoiding
all of them, within the strict contour and accuracy margins.  Different
principal parts then separate by a divergent power in real part.
The terms $O(\log B)$ in \eqref{phase:expansion-eq} cannot offset this
separation.  Hence the dominant end contributions form one class with
identical $P_j$.

Choose a representative $j_0$ of this class.  Identical principal parts
imply identical leading $c_j$, and thus
$Z_{0,j}/Z_{0,j_0}\to1$ on the same branches.
By \eqref{phase:expansion-eq}, \eqref{phase:real-ratio}, and the real
Laurent constants,
\[
 Z_{*,j}^2-Z_{*,j_0}^2\longrightarrow d_j\in\mathbb R,
       \qquad \frac{Z_{*,j}^{-2}}{Z_{*,j_0}^{-2}}\longrightarrow1.
\]
It follows that
\begin{equation}
 \frac{Z_{*,j}^{-2}e^{-Z_{*,j}^2/4}}
      {Z_{*,j_0}^{-2}e^{-Z_{*,j_0}^2/4}}
       \longrightarrow e^{-d_j/4}>0.
 \label{phase:positive-ratios}
\end{equation}
Thus no cancellation is possible within the dominant class.  Ends in
other classes have ratios tending to zero exponentially in a positive
power of $\rho^{-1}$.

The sign and phase of the coefficient multiplying each contribution
also agree.  Proposition~\ref{jump:jump} uses the branch of the area
element continued from positive real area and the outward integration
convention on each end; its boundary momentum calculation gives the same
nonzero constant for the same order of the two lateral bypasses.
For ends with a tied dominant principal part the opening pole is crossed
on the same side, so the two labels are ordered alike.  An ambient axis
orientation or a choice of normal cannot reverse one term independently.
Consequently that proposition and \eqref{phase:positive-ratios} give
\begin{equation}
 G_+-G_-=C_{\rm St}
    Z_{*,j_0}^{-2}e^{-Z_{*,j_0}^2/4}
       \left(\sum_{j\ \mathrm{dominant}}e^{-d_j/4}+o(1)\right),
 \qquad C_{\rm St}\ne0.
 \label{phase:nonzero-jump}
\end{equation}
Here the error is justified in the specified order: choose the exterior
translator cutoff $Y_0$ sufficiently large that its relative error is
smaller than a fixed fraction of the positive sum, and then take $B$
sufficiently large.  Equivalently the two-limit error in
Proposition~\ref{jump:jump} yields the lower bound in
\eqref{phase:nonzero-jump}; an interchange of the two limits is unnecessary.

Finally retain the leading action scale from \eqref{arr:A0}, namely
\[
 A_{\mathrm{lead}}(s)=\min_{\mathrm{adverse}\ j}
       \frac{|(-c_js^k)^{-1/2}|^2}{4}.
\]
The actual poles in \eqref{phase:expansion-eq} satisfy
\[
 \frac{Z_{0,j}^2}{(-c_js^k)^{-1}}=1+o(1)
\]
uniformly on the finite adverse list.  After the fixed generic angle
is chosen sufficiently close to the initial ray, one adverse end $j_1$
satisfies
\[
 \frac{\operatorname{Re}(-c_{j_1}s^k)^{-1}}4
       \le 1.05A_{\mathrm{lead}}.
\]
By the choice of the dominant class,
$\operatorname{Re}Z_{*,j_0}^2\le\operatorname{Re}Z_{*,j_1}^2+o(B^2)$.
The pole comparison above and \eqref{phase:expansion-eq} therefore give
\[
 \frac{\operatorname{Re}Z_{*,j_0}^2}{4}
             \le(1.05+o(1))A_{\mathrm{lead}}.
\]
Moreover $|Z_{*,j_0}|\asymp B$, and the positive sum in
\eqref{phase:nonzero-jump} is bounded away from zero.  Its fixed nonzero
coefficient and inverse-square prefactor change the logarithm of the
lower bound by only $O(\log B)=o(A_{\mathrm{lead}})$.  For all
sufficiently small $s$ on the chosen ray, that lower bound is consequently
\[
       |G_+-G_-|\ge e^{-1.10A_{\mathrm{lead}}}.
\]
This is incompatible with
\[
       |G_+-G_-|=O(e^{-1.2A_{\mathrm{lead}}})
\]
from Proposition~\ref{arr:action-upper}.
The contradiction excludes the assumed arc.
\end{proof}
